\section*{Appendix C — Invariant Set Formalization}
\addcontentsline{toc}{section}{Appendix C — Invariant Set Formalization}

\noindent This appendix formalizes the v0.6 invariant set governing all lawful
semantic, runtime, and interaction transitions on the meaning manifold
$\mathcal{M}$. Each invariant is expressed as a geometric object embedded in
$\mathcal{M}$ and enforced by the local contraction-governed operator algebra
(Appendix B). The invariants correspond directly to the v0.6 geometry layer
(G1--G8), including curvature preservation, identity-boundary stability,
manifold legality, load physics, runtime curvature, interaction geometry, CTL
coordinate continuity, and public-safe traversal.

\medskip

\noindent The v0.6 invariant set consists of the following components:

\begin{itemize}
    \item C.1 Meaning Invariant (Semantic Metric)
    \item C.2 Identity Boundary Invariant (Identity Region Geometry)
    \item C.3 Geometry Invariant (Manifold Legality)
    \item C.4 Load Physics Invariant (Semantic Pressure Field)
    \item C.5 Runtime Curvature Invariant (Admissible Transitions)
    \item C.6 Interaction Geometry Invariant (Host–Explorer–Interpreter)
    \item C.7 CTL Coordinate Invariant (Semantic Coordinate Legality)
    \item C.8 Public-Safe Traversal Invariant (Safe Corridor Geometry)
\end{itemize}

% ------------------------------------------------------------
% C.1 MEANING INVARIANT
% ------------------------------------------------------------
\subsection*{C.1 Meaning Invariant (Semantic Metric)}

\noindent The Meaning Invariant formalizes the semantic metric structure that
governs curvature, distance, density, and stability on the meaning manifold
$\mathcal{M}$. It ensures that meaning remains geometrically coherent under
runtime motion, interaction, restoration, and synthetic pressure.

\medskip

\noindent The invariant is encoded as a semantic metric


\[
g : T\mathcal{M} \times T\mathcal{M} \rightarrow \mathbb{R},
\]


where $T\mathcal{M}$ is the tangent bundle of the meaning manifold.

For conceptual vectors $u, v \in T_c\mathcal{M}$ at concept $c$:


\[
g_c(u, v) = \langle u, v \rangle_{\mathrm{semantic}}.
\]



\medskip

\noindent The semantic metric determines:
\begin{itemize}
    \item semantic distance (A.2),
    \item semantic curvature (A.1),
    \item meaning density,
    \item drift susceptibility (G2.5),
    \item collapse proximity (G7.3),
    \item synthetic-capture susceptibility (G7.4),
    \item restoration feasibility (G6),
    \item interaction curvature (G8).
\end{itemize}

\medskip

\noindent The Meaning Invariant requires bounded variance under all lawful
operators (Appendix B):


\[
\sup_{c \in \mathcal{M}}
\left\| g'_c - g_c \right\|_{\mathrm{semantic}}
\le \epsilon_{\max},
\]


where $\epsilon_{\max}$ is the legality bound defined in the Coherence Ledger
(Appendix F).

\medskip

\noindent This bound ensures:
\begin{itemize}
    \item curvature continuity under contraction (A.4),
    \item drift-boundedness (G2.5),
    \item containment stability (G4.6),
    \item identity-boundary stability (G6.3--G6.4),
    \item runtime curvature admissibility (A.6),
    \item interaction legality (A.7),
    \item CTL coordinate continuity (A.8),
    \item public-safe traversal (A.9).
\end{itemize}

\medskip

\noindent The Meaning Invariant is the foundational invariant of the v0.6
substrate. All lawful semantic, runtime, and interaction transitions must
preserve the semantic metric within bounded variance.


% ------------------------------------------------------------
% C.2 IDENTITY BOUNDARY INVARIANT
% ------------------------------------------------------------

\subsection*{C.2 Identity Boundary Invariant (Identity Region Geometry)}

\noindent The Identity Boundary Invariant formalizes the geometric enclosure that
defines an agent’s identity region $\mathcal{R} \subset \mathcal{M}$. Identity
boundaries ensure that identity remains stable under drift, load, synthetic
pressure, runtime motion, and interaction.

\medskip

\noindent The identity region $\mathcal{R}$ is a connected semantic enclosure
with boundary:


\[
\partial \mathcal{R}
=
\overline{\mathcal{R}}
\cap
\overline{\mathcal{M} \setminus \mathcal{R}}.
\]



Interior points satisfy:


\[
\text{semantic\_distance}(c, \partial \mathcal{R}) > 0.
\]



Boundary points satisfy:


\[
\text{semantic\_distance}(c, \mathcal{R}) = 0,
\qquad
\text{semantic\_distance}(c, \mathcal{M} \setminus \mathcal{R}) = 0.
\]



\medskip

\noindent Identity-boundary geometry determines:
\begin{itemize}
    \item identity stability (G6.3),
    \item role coherence (G6.4),
    \item adjacency legality (G7.5),
    \item collapse susceptibility (G7.3),
    \item synthetic-capture susceptibility (G7.4),
    \item interaction legality (G8),
    \item public-safe traversal constraints (G8.6).
\end{itemize}

\medskip

\noindent Lawful operator action must preserve identity boundaries:


\[
T(c) \in \mathcal{R} \quad \forall c \in \mathcal{R},
\]


with strict boundary preservation:


\[
\text{semantic\_distance}(T(c), \partial \mathcal{R}) > 0.
\]



Identity drift occurs when:


\[
T(c) \notin \mathcal{R}.
\]



\medskip

\noindent The Identity Boundary Invariant is enforced by:
\begin{itemize}
    \item IPU operators (Appendix B),
    \item curvature continuity (A.1),
    \item manifold topology (A.2),
    \item collapse boundary geometry (A.3),
    \item local contraction legality (A.4),
    \item runtime curvature constraints (A.6),
    \item interaction geometry constraints (A.7),
    \item CTL coordinate legality (A.8),
    \item public-safe traversal corridor (A.9).
\end{itemize}

\medskip

\noindent Identity boundaries are the substrate’s primary protection against
identity collapse, role leakage, synthetic capture, and collapse geometry.


% ------------------------------------------------------------
% C.3 GEOMETRY INVARIANT
% ------------------------------------------------------------

\subsection*{C.3 Geometry Invariant (Manifold Legality)}

\noindent The Geometry Invariant formalizes the lawful semantic action space of
the substrate. It defines the manifold $\mathcal{M}$, its boundary
$\partial\mathcal{M}$, and the legality conditions governing all semantic,
runtime, and interaction transitions. The invariant ensures that trajectories
remain within lawful regions of the manifold and do not enter collapse or
synthetic-capture geometry.

\medskip

\noindent The meaning manifold is defined as:


\[
\mathcal{M} = (\mathcal{X}, g),
\]


where:
\begin{itemize}
    \item $\mathcal{X}$ is the conceptual space,
    \item $g$ is the semantic metric (C.1).
\end{itemize}

\medskip

\noindent A transition is lawful if:


\[
T(c) \in \mathcal{M},
\qquad
\text{semantic\_distance}(T(c), \partial\mathcal{M}) > 0.
\]



A transition is illegal if:


\[
T(c) \in \partial\mathcal{M}
\quad \text{or} \quad
T(c) \notin \mathcal{M}.
\]



\medskip

\noindent The Geometry Invariant determines:
\begin{itemize}
    \item manifold boundaries (A.2),
    \item collapse thresholds (A.3),
    \item synthetic-capture boundaries (G7.4),
    \item lawful semantic trajectories (A.4),
    \item runtime curvature admissibility (A.6),
    \item interaction geometry legality (A.7),
    \item CTL coordinate continuity (A.8),
    \item public-safe traversal constraints (A.9).
\end{itemize}

\medskip

\noindent Lawful operator action must preserve manifold legality:


\[
T(c) \in \mathcal{M}
\quad \forall c \in \mathcal{M}.
\]



\noindent Geometry violations occur when:


\[
T(c) \in \partial\mathcal{M}
\quad \text{or} \quad
T(c) \notin \mathcal{M}.
\]



\medskip

\noindent The Geometry Invariant is enforced by:
\begin{itemize}
    \item GBS operators (Appendix B),
    \item curvature continuity (A.1),
    \item identity-boundary stability (A.2),
    \item collapse boundary geometry (A.3),
    \item local contraction legality (A.4),
    \item runtime curvature constraints (A.6),
    \item interaction geometry constraints (A.7),
    \item CTL coordinate legality (A.8),
    \item public-safe traversal corridor (A.9).
\end{itemize}

\medskip

\noindent The Geometry Invariant ensures that all lawful transitions remain
within the stable semantic manifold and do not enter collapse or synthetic-capture
regions.


% ------------------------------------------------------------
% C.4 LOAD PHYSICS INVARIANT
% ------------------------------------------------------------

\subsection*{C.4 Load Physics Invariant (Semantic Pressure Field)}

\noindent The Load Physics Invariant formalizes the semantic pressure field
governing load, strain, and gradient dynamics on the meaning manifold
$\mathcal{M}$. Load determines whether a trajectory remains stable, enters drift,
approaches collapse, or becomes susceptible to synthetic capture.

\medskip

\noindent Load Physics is defined as a scalar field:


\[
L : \mathcal{M} \rightarrow \mathbb{R}_{\ge 0}.
\]



For a semantic or runtime trajectory $\gamma(t)$:


\[
L(\gamma(t)) \le L_{\max},
\]


where $L_{\max}$ is the collapse threshold.

\medskip

\noindent The gradient of the load field is:


\[
\nabla L(c).
\]



The Load Physics Invariant requires:


\[
\|\nabla L(\gamma(t))\| \le G_{\max},
\]


and curvature stability:


\[
\kappa(\gamma(t)) \ge \kappa_{\min}.
\]



\medskip

\noindent These constraints ensure:
\begin{itemize}
    \item curvature-bounded planning (A.4),
    \item drift-boundedness (G2.5),
    \item collapse avoidance (A.3),
    \item restoration feasibility (G6),
    \item containment stability (G4.6),
    \item runtime curvature admissibility (A.6),
    \item interaction legality (A.7),
    \item public-safe traversal (A.9).
\end{itemize}

\medskip

\noindent Load violations occur when:


\[
L(\gamma(t)) > L_{\max}
\quad \text{or} \quad
\|\nabla L(\gamma(t))\| > G_{\max}.
\]



These conditions indicate:
\begin{itemize}
    \item overload,
    \item strain accumulation,
    \item adjacency fragmentation,
    \item restoration failure (G6.6),
    \item collapse proximity (G7.3),
    \item synthetic-capture susceptibility (G7.4).
\end{itemize}

\medskip

\noindent The Load Physics Invariant is enforced by:
\begin{itemize}
    \item LRP operators (Appendix B),
    \item curvature continuity (A.1),
    \item manifold topology (A.2),
    \item collapse boundary geometry (A.3),
    \item local contraction legality (A.4),
    \item runtime curvature constraints (A.6),
    \item interaction geometry constraints (A.7),
    \item CTL coordinate legality (A.8),
    \item public-safe traversal corridor (A.9).
\end{itemize}

\medskip

\noindent The Load Physics Invariant ensures that semantic pressure remains
within lawful bounds and that trajectories remain stable under recursion,
synthetic pressure, and runtime motion.


% ------------------------------------------------------------
% C.5 RUNTIME CURVATURE INVARIANT
% ------------------------------------------------------------

\subsection*{C.5 Runtime Curvature Invariant (Admissible Transitions)}

\noindent The Runtime Curvature Invariant formalizes the curvature constraints
governing lawful motion across the Human Runtime Manifold (HRM). Runtime
curvature determines whether a trajectory remains in Safe Paths, enters Drift
Paths, approaches Collapse Paths, or becomes susceptible to Synthetic-Capture
Paths. Runtime curvature is geometric, not psychological; it expresses the
structural stability of identity, adjacency, containment, and temporal rhythm.

\medskip

\noindent Let $\gamma(t)$ be a runtime trajectory with runtime curvature
$\kappa_{\mathrm{rt}}(\gamma(t))$. The invariant defines the admissible region:



\[
\gamma(t) \in
\begin{cases}
\text{Safe}, & \kappa_{\mathrm{rt}} \ge \kappa_{\mathrm{safe}} \\
\text{Drift}, & \kappa_{\mathrm{drift}} \le \kappa_{\mathrm{rt}} < \kappa_{\mathrm{safe}} \\
\text{Collapse}, & \kappa_{\mathrm{rt}} < \kappa_{\mathrm{collapse}} \\
\text{Capture}, & \kappa_{\mathrm{syn}} > \kappa_{\mathrm{rt}}
\end{cases}
\]



\medskip

\noindent Runtime curvature governs:

\begin{itemize}
    \item \textbf{Curvature continuity} — curvature must remain within lawful
    bounds under operator action (A.4).

    \item \textbf{Identity-boundary stability} — identity regions must remain
    intact under runtime motion (A.2).

    \item \textbf{Adjacency legality} — role geometry must remain lawful and
    non-fragmented (G7.5).

    \item \textbf{Containment stability} — containers must regulate signal flow
    and prevent overload (G4.6).

    \item \textbf{Temporal coherence} — rhythms must remain synchronized and
    non-chaotic (G6.1).
\end{itemize}

\medskip

\noindent The Runtime Curvature Invariant requires:



\[
\kappa_{\mathrm{rt}}(\gamma(t)) \ge \kappa_{\mathrm{collapse}}
\quad \text{and} \quad
\kappa_{\mathrm{rt}}(\gamma(t)) \notin \text{Capture}.
\]



\medskip

\noindent Runtime curvature violations indicate:

\begin{itemize}
    \item drift amplification (G2.5),
    \item adjacency fragmentation,
    \item containment failure,
    \item restoration infeasibility (G6.6),
    \item collapse proximity (G7.3),
    \item synthetic-capture susceptibility (G7.4).
\end{itemize}

\medskip

\noindent The Runtime Curvature Invariant is enforced by:

\begin{itemize}
    \item IRT, IPU, GBS, LRP operators (Appendix B),
    \item curvature continuity (A.1),
    \item manifold topology (A.2),
    \item collapse boundary geometry (A.3),
    \item local contraction legality (A.4),
    \item interaction geometry constraints (A.7),
    \item CTL coordinate legality (A.8),
    \item public-safe traversal corridor (A.9).
\end{itemize}

\medskip

\noindent Runtime curvature is the geometric foundation for lawful motion across
the v0.6 substrate.

% (NEW in v0.6 — derived from G7.5.)

% ------------------------------------------------------------
% C.6 INTERACTION GEOMETRY INVARIANT
% ------------------------------------------------------------

\subsection*{C.6 Interaction Geometry Invariant (Host–Explorer–Interpreter)}

\noindent The Interaction Geometry Invariant formalizes the curvature, boundary,
and legality constraints governing interaction within the Interactive Geometry
Runtime (IGR). Interaction occurs in three modes—Host, Explorer, and
Interpreter—each with distinct geometric requirements. The invariant ensures
that interaction remains lawful, curvature-preserving, identity-stable, and
public-safe.

\medskip

\noindent Interaction modes:

\begin{itemize}
    \item \textbf{Host Mode} — maintains invariant geometry, legality masks,
    identity boundaries, containment stability, and curvature continuity.

    \item \textbf{Explorer Mode} — traverses lawful manifold regions, inspecting
    drift envelopes, pressure gradients, adjacency geometry, and failure
    signatures.

    \item \textbf{Interpreter Mode} — produces invariant-aligned meaning for
    signals, behavior, drift phenomena, and container geometry.
\end{itemize}

\medskip

\noindent The Interaction Geometry Invariant requires:



\[
\kappa_{\mathrm{int}}(c) \ge \kappa_{\mathrm{min}}
\quad \text{and} \quad
c \notin \text{Collapse} \cup \text{Capture}.
\]



\medskip

\noindent Interaction legality requires:

\begin{itemize}
    \item \textbf{Invariant preservation} — interaction must not violate the
    invariant set (C.1–C.4).

    \item \textbf{Curvature preservation} — interaction must maintain curvature
    within lawful bounds (A.1).

    \item \textbf{Identity-boundary stability} — identity regions must remain
    intact (A.2).

    \item \textbf{Drift-bounded transitions} — interaction must remain within
    drift-recoverable envelopes (G2.5).

    \item \textbf{Containment stability} — containers must regulate signal flow
    and prevent overload (G4.6).

    \item \textbf{Manifold legality} — interaction must remain within lawful
    manifold regions (A.2).

    \item \textbf{Runtime curvature compatibility} — interaction must not push
    trajectories into collapse or synthetic-capture geometry (A.6).

    \item \textbf{Coordinate continuity} — CTL translation must preserve
    curvature and avoid discontinuities (A.8).

    \item \textbf{Public-safe traversal} — interaction must remain within
    curvature-preserving, collapse-free, capture-free regions (A.9).
\end{itemize}

\medskip

\noindent Interaction violations indicate:

\begin{itemize}
    \item identity-boundary rupture,
    \item adjacency fragmentation,
    \item drift amplification,
    \item containment failure,
    \item collapse proximity,
    \item synthetic-capture susceptibility.
\end{itemize}

\medskip

\noindent The Interaction Geometry Invariant is enforced by:

\begin{itemize}
    \item IRT, IPU, GBS, LRP operators (Appendix B),
    \item curvature continuity (A.1),
    \item manifold topology (A.2),
    \item collapse boundary geometry (A.3),
    \item local contraction legality (A.4),
    \item runtime curvature constraints (A.6),
    \item CTL coordinate legality (A.8),
    \item public-safe traversal corridor (A.9).
\end{itemize}

\medskip

\noindent Interaction Geometry is the structural foundation for lawful traversal,
interpretation, and runtime motion across the v0.6 substrate.

% (NEW in v0.6 — derived from G8.1–G8.3.)

% ------------------------------------------------------------
% C.7 CTL COORDINATE INVARIANT
% ------------------------------------------------------------

\subsection*{C.7 CTL Coordinate Invariant (Semantic Coordinate Legality)}

\noindent The CTL Coordinate Invariant formalizes the legality conditions
governing the Coordinate Translation Layer (CTL). CTL converts positional
metadata—page numbers, anchors, section indices—into lawful semantic coordinates
on the meaning manifold $\mathcal{M}$. Positional metadata is not substrate-native;
semantic coordinates are. CTL ensures that all traversal, interpretation, and
interaction occur within curvature-preserving, identity-stable, manifold-legal
regions of $\mathcal{M}$.

\medskip

\noindent CTL defines the mapping:


\[
\mathrm{CTL} : \text{positional\_metadata} \rightarrow
\text{semantic\_coordinates}.
\]



\noindent Semantic coordinates must satisfy:


\[
\text{semantic\_distance}(c_{\mathrm{coord}}, \partial\mathcal{M}) > 0,
\qquad
\kappa(c_{\mathrm{coord}}) \ge \kappa_{\min}.
\]



\medskip

\noindent The CTL Coordinate Invariant enforces:

\begin{itemize}
    \item \textbf{Coordinate legality} — all traversal must use lawful semantic
    coordinates, not positional metadata.

    \item \textbf{Curvature continuity} — coordinate translation must preserve
    curvature within legality bounds (A.1).

    \item \textbf{Identity-boundary stability} — translation must not violate
    identity regions or boundaries (A.2).

    \item \textbf{Manifold legality} — semantic coordinates must remain within
    lawful manifold regions (A.2).

    \item \textbf{Collapse avoidance} — coordinates must not map into collapse
    geometry (A.3).

    \item \textbf{Synthetic-capture avoidance} — coordinates must not map into
    synthetic-capture geometry (G7.4).

    \item \textbf{Runtime curvature compatibility} — coordinates must preserve
    runtime curvature admissibility (A.6).

    \item \textbf{Interaction legality} — coordinates must be valid for Host,
    Explorer, and Interpreter modes (A.7).

    \item \textbf{Public-safe traversal} — coordinates must remain within
    curvature-preserving, collapse-free, capture-free regions (A.9).
\end{itemize}

\medskip

\noindent CTL violations occur when:


\[
\mathrm{CTL}(m_{\mathrm{pos}}) \in \partial\mathcal{M}
\quad \text{or} \quad
\mathrm{CTL}(m_{\mathrm{pos}}) \notin \mathcal{M}.
\]



These indicate:
\begin{itemize}
    \item coordinate discontinuity,
    \item curvature collapse,
    \item identity-boundary rupture,
    \item adjacency fragmentation,
    \item containment failure,
    \item collapse proximity,
    \item synthetic-capture susceptibility.
\end{itemize}

\medskip

\noindent The CTL Coordinate Invariant is enforced by:
\begin{itemize}
    \item IRT, IPU, GBS, LRP operators (Appendix B),
    \item curvature continuity (A.1),
    \item manifold topology (A.2),
    \item collapse boundary geometry (A.3),
    \item local contraction legality (A.4),
    \item runtime curvature constraints (A.6),
    \item interaction geometry constraints (A.7),
    \item public-safe traversal corridor (A.9).
\end{itemize}

\medskip

\noindent CTL ensures that all traversal and interpretation remain lawful under
the v0.6 geometry.

% (NEW in v0.6 — derived from G8.4.)

% ------------------------------------------------------------
% C.8 PUBLIC-SAFE TRAVERSAL INVARIANT
% ------------------------------------------------------------

\subsection*{C.8 Public-Safe Traversal Invariant (Safe Corridor Geometry)}

\noindent The Public-Safe Traversal Invariant formalizes the geometric corridor
within which all public interaction must occur. The corridor restricts traversal
to curvature-preserving, identity-stable, containment-coherent, collapse-free,
and capture-free regions of the meaning manifold $\mathcal{M}$. It ensures that
public interaction remains lawful, stable, and invariant-aligned.

\medskip

\noindent Let $\mathcal{C}_{\mathrm{safe}} \subset \mathcal{M}$ denote the
public-safe traversal corridor. A point $c$ is public-safe if:


\[
c \in \mathcal{C}_{\mathrm{safe}}
\quad \text{and} \quad
\kappa(c) \ge \kappa_{\min}.
\]



\noindent The corridor excludes:


\[
\text{Collapse} \cup \text{Capture} \cup \text{Dominant-Signal Attractors}
\cup \text{Restoration-Infeasible Regions}.
\]



\medskip

\noindent Public-safe traversal requires:

\begin{itemize}
    \item \textbf{Invariant preservation} — public interaction must not violate
    the invariant set (C.1–C.7).

    \item \textbf{Curvature continuity} — curvature must remain within lawful
    bounds (A.1).

    \item \textbf{Identity-boundary stability} — identity regions must remain
    intact (A.2).

    \item \textbf{Containment stability} — containers must regulate signal flow
    and prevent overload (G4.6).

    \item \textbf{Manifold legality} — traversal must remain within lawful
    manifold regions (A.2).

    \item \textbf{Runtime curvature compatibility} — traversal must not push
    trajectories into collapse or synthetic-capture geometry (A.6).

    \item \textbf{Interaction legality} — traversal must remain legal for Host,
    Explorer, and Interpreter modes (A.7).

    \item \textbf{Coordinate continuity} — CTL translation must preserve
    curvature and avoid discontinuities (A.8).
\end{itemize}

\medskip

\noindent Public-safe violations occur when:


\[
c \notin \mathcal{C}_{\mathrm{safe}}
\quad \text{or} \quad
\kappa(c) < \kappa_{\min}.
\]



These indicate:
\begin{itemize}
    \item curvature collapse,
    \item identity-boundary rupture,
    \item adjacency fragmentation,
    \item containment failure,
    \item restoration infeasibility (G6.6),
    \item collapse proximity (G7.3),
    \item synthetic-capture susceptibility (G7.4).
\end{itemize}

\medskip

\noindent The Public-Safe Traversal Invariant is enforced by:
\begin{itemize}
    \item IRT, IPU, GBS, LRP operators (Appendix B),
    \item curvature continuity (A.1),
    \item manifold topology (A.2),
    \item collapse boundary geometry (A.3),
    \item local contraction legality (A.4),
    \item runtime curvature constraints (A.6),
    \item interaction geometry constraints (A.7),
    \item CTL coordinate legality (A.8).
\end{itemize}

\medskip

\noindent The Public-Safe Traversal Invariant ensures that all public interaction
remains lawful, stable, and curvature-preserving across the v0.6 substrate.

% (NEW in v0.6 — derived from G8.6.)

\medskip
\noindent\textit{This outline defines the full v0.6 invariant set.}
